\section{系统概述}\label{sec:System}

\subsection{系统的基本概念}\label{sub:concepts_of_system}

对于单输入单输出系统，将输入信号称为\textbf{激励} (excitation)，输出信号称为\textbf{响应} (response)，并将时域信号分别用$e(t),r(t)$表示，如果系统用字母H表示，可以记$r(t)=H[e(t)]$，或更简洁地，$r(t)=He(t)$。$H$是函数空间上的映射。

\begin{example}
    给定以下一阶动态电路，并设各元件无初始储能。
    \begin{center}
        \begin{circuitikz}[american]
            \draw (0,0) to[V=$1$V] (0,3)
            to[closing switch] (5,3)
            to[C=$C$] (5,0)
            to[R=$R$] (0,0);
        \end{circuitikz}
    \end{center}
    在$t=0$时刻闭合开关，则系统的激励为$e(t)=u(t)$。如果认为电容两端的电压$u_C(t)$是系统的响应$r(t)$，则可以用$r(t)=H[e(t)]$来表示这个系统。

    使用三要素法分析该电路。初状态下$u_C(0)=0$V，末状态下$u_C(\infty)=1$V，时间常数$\tau=RC$，因此
    \[e(t)=u_C(t)=1-\mathrm{e}^{t/RC}\]
\end{example}

本章主要介绍连续系统，即激励与响应均为连续信号的系统。与之相对，激励与响应均为离散信号的系统称为离散系统，我们将在第六章和第七章专门讨论离散系统。

一般来说，如果我们对连续系统进行适当的建模，则可以用函数方程来描述这个系统。函数方程不意味着我们一定能够给出响应关于激励的封闭表达式。
\begin{example}
设激励$e(t)$与响应$r(t)$满足
\[r^3(t)+4r(t)=\exp e(t)\]
这是一个超越方程，但由于$f(x)=x^3+4x$严格单调递增，可由激励唯一确定响应的值，这也是一个系统。
\end{example}

其中，我们特别关注能够用\textbf{常系数线性微分方程（组）}描述的系统，例如，第一个例子中的动态电路就可以用一阶常系数线性微分方程来描述。我们首先给出微分方程的概念：
\begin{definition}[微分方程]
    称如下形式的方程为\textbf{常微分方程}(ordinary differential equation,ODE)\cite{ode}：
    \[F(x,y,y',\cdots,y^{(n)})=0,n\geq 1\]
    其中$F$为已知函数，自变量$x$在$\mathbb{R}$的某个区间上变化，未知函数$y\in\mathbb{R}$，正整数$n$称为方程的\textbf{阶}。如果$n\geq 2$，则方程称为\textbf{高阶微分方程}。
\end{definition}
\begin{definition}[初值问题]
    数学上，将微分方程中的函数在某一点的各阶导数值称为\textbf{初始条件}，以$t=0$时刻为例，又称为\textbf{0时刻状态}：
    \[y(0)=y_0,y'(0)=y'_0,\cdots,y^{(n-1)}(0)=y^{(n-1)}_0\]
    我们称这样的问题为\textbf{初值问题}。
\end{definition}

常微分方程的解空间是有限维的，只要给定足够的初值数量，就可以唯一确定方程的解，因此常微分方程系统都具有\textbf{即时性}。在\ref{sub:compare_to_high_order}中，我们将对线性微分方程组的解空间维数进行说明。

\begin{definition}[即时性与记忆性]
一个系统具有\textbf{即时性}（或无记忆性、Markov性），如果该系统在任意时刻$t$的状态完全且唯一地取决于该时刻的系统状态与激励，而不显式地依赖于任何过去时刻$t'<t$的系统状态或历史激励。换言之，如果一个初值问题的解存在且唯一，则由它描述的系统具有即时性。

一个系统具有\textbf{记忆性}（或历史依赖性、非Markov性），如果系统的响应在任意时刻$t$的状态演化率显式地依赖于过去时刻$t'<t$的系统状态或历史激励，而不能被简化为仅依赖于当前时刻$t$的量。
\end{definition}

\begin{example}[积分方程]
    设系统由积分方程$r(t)=\int_{-\infty}^{t}e(y)\,dy$描述，则它是记忆系统。
\end{example}
\begin{example}[时滞系统]
    由时滞微分方程描述的系统称为时滞系统。本书不引入这个概念，仅仅给出以下例子：
    \[y'(t)=y(t)+y(t-\tau)\]
    其中$\tau$是常数。一般来说，我们不将时滞微分方程归类为微分方程。时滞微分方程的解空间是无限维的，不可能通过某一时刻的系统状态唯一确定方程的解，因为我们只能指出有限个系统状态。这个方程描述了一个无激励的记忆系统。
\end{example}

下面介绍一些常见的系统特性。

\begin{definition}[线性系统]
    如果算子H是线性的，即
\[\forall a,b\in\mathbb{C},\forall e_1(t),e_2(t),H[ae_1(t)+be_2(t)]=aH[e_1(t)]+bH[e_2(t)]\]
就说系统是\textbf{线性系统}(linear system)，且满足叠加法则 (principle of superposition)。
\end{definition}

由于系统原有的能量，即便激励$e(t)=0$，响应$r(t)$也可能不为零。对于线性系统，可以将$r(t)$分解为两部分：
\begin{definition}[零输入响应与零状态响应]
   \textbf{零输入响应}(zero input response,$r_{zi}(t)$)是指$e(t)=0$时，由于系统原有能量而产生的响应；\textbf{零状态响应}(zero state response,$r_{zs}(t)$)是指将系统初始能量置零时的响应。\textbf{全响应}$r(t)=r_{zi}(t)+r_{zs}(t)$。
\end{definition}
对于零输入响应不为0的线性系统，考虑以两个激励$e_1(t),e_2(t)$的线性组合作为输入，则按照先前对线性系统的定义，直接将对应的两个响应$r_1(t),r_2(t)$进行线性组合，则零输入响应被重复计算，这不是我们希望看到的，因此我们其实应该定义\textbf{零状态响应有线性性的系统为线性系统}。后文中，如无特殊说明，响应总应理解为零状态响应，而零输入响应仅仅在分析给定系统时单独研究。

\begin{definition}[时不变性和移不变性]
    如果一个系统的输出不依赖于输入信号施加于系统的时间，输入信号发生时移，输出信号也发生相同的时移，即
    $$\forall b\in\mathbb{R},H[e(t-b)]=r(t-b),\text{即}H[\tau_b e(t)]=\tau_b r(t)$$
    则称该系统是\textbf{时不变的}(time-invariant)。如果激励和响应的自变量是空间，我们同样可以定义系统的\textbf{移不变性}(shift-invariance)，这时$b$是空间中的向量。后文中一般不特意区分时域和空间域。
\end{definition}
\begin{definition}[线性时不变系统]
    如果一个系统既是线性的又是时不变的，就说它是\textbf{线性时不变系统}(linear time-invariant system,LTI)。
\end{definition}

线性时不变系统是信号与系统课程中所研究的主要系统类型。下面两节中我们将看到，常系数线性微分方程（组）描述的系统是线性时不变系统。对于线性时不变系统，定义：

\begin{definition}[单位脉冲响应与单位阶跃响应]
    当$e(t)=\delta(t)$时，零状态响应$r_{zs}(t)$为\textbf{单位脉冲响应}；当$e(t)=u(t)$时，零状态响应$r_{zs}(t)$称为\textbf{单位阶跃响应}。习惯上，将单位脉冲响应记为$h(t)$，单位阶跃响应记为$g(t)$。
\end{definition}

除了线性性与时不变性，系统还有一些其他的特性，例如稳定性和可逆性。系统的其他特性将在后文中陆续介绍。

\begin{definition}[稳定性]
如果一个系统在激励信号有界时，响应也是有界的，则称之为\textbf{稳定系统}(bounded-input bounded-output,BIBO)，即
\[e(t)\in L^1(\mathbb{R})\implies r(t)\in L^1(\mathbb{R})\]
\end{definition}
我们将在\ref{sub:stability_of_system}和\ref{sub:discrete_system_function}中，分别借助拉普拉斯变换和z变换的理论，给出连续的和离散的线性时不变系统具有稳定性的充分必要条件。

\begin{definition}[可逆性]
    如果H（作为映射）如果是单射，则也是双射（因为我们不关注其值域），于是H是可逆的。换句话说，一种响应仅可能对应唯一的激励。
\end{definition}

\subsection{常见线性系统}\label{sub:ordinary_system}

\begin{example}[时域乘积]
\begin{equation*}
    r(t)=f(t)e(t)
\end{equation*}
例如描述开关开闭的系统，$f(t)=\chi_{[T_1,T_2]}(t)$或$u(t)$。
\end{example}

\begin{example}[微分器与积分器]
    积分器是由积分方程$$r(t)=\int_{-\infty}^{t}e(y)\,dy$$描述的系统；微分器是由微分方程$$r(t)=\frac{de(t)}{dt}$$描述的系统。它们可以通过运算放大器实现。
\end{example}

\begin{example}[对核积分]
\begin{equation*}
    r(t)=\int_{a}^{b}k(x,y)e(y)\,dy
\end{equation*}
其中，$k(x,y)$称为\textbf{核函数}(kernel)，$\int_{a}^{b}k(x,y)e(y)\,dy$称为\textbf{对核函数积分}(integrating against a kernel)，简称\textbf{对核积分}。
\end{example}
根据积分的线性性，对核积分系统也是线性的。习惯上将核函数的符号取为$k$。我们将在下一小节中看到，线性系统总能描述为对核积分系统，因此它是我们考虑的主要的连续线性系统的类型。利用核函数可以定义：

\begin{definition}[转置系统对称系统与自伴系统]
    类比线性代数中算子的转置以及对称算子、自伴算子，对于对核积分系统
    $r(t)=\int_{a}^{b}k(x,y)e(y)\,dy$
    有：
\begin{itemize}
    \item 转置系统$H^T$描述为$H^T[e(t)]=\int_{a}^{b}k(y,x)e(y)\,dy$
    \item 如果$k(x,y)=k(y,x)$，称k是\textbf{对称的}并将系统称为\textbf{对称系统}
    \item 如果$k(x,y)=\overline{k(y,x)}$，称k是\textbf{埃尔米特的}并将系统称为\textbf{自伴系统}。
\end{itemize}
\end{definition}

%类似地，对于周期离散信号，矩阵乘法定义了线性系统：
%\begin{example}[矩阵乘法]
%设离散系统的激励与响应均为$n$维离散信号，并且响应为激励的线性组合，则\begin{equation*}
%    \mathbf{w}=A\mathbf{v},A\in M_{n\times n}(\mathbb{C})
%\end{equation*}
%\end{example}
%线性代数中，如果一个线性算子L的在某一组基下的矩阵表示是A，则将L的转置定义为用$A^T$表示的线性算子，其中T表示矩阵的转置，如果矩阵A是对称的，则算子L称为对称算子；L的共轭转置定义为用$A^H$表示的线性算子，其中H表示矩阵的共轭转置，如果A是埃尔米特的(Hermitian)，即$A^H=A$，则称算子L是埃尔米特的，并说L是\textbf{自伴算子}(self-adjoint operator).对于矩阵乘法系统，可以根据A的性质定义\textbf{对称系统、自伴系统}。

\begin{example}[卷积]
    取定衰减足够快的连续信号$h$，$r(t)=h*e(t)$.
\end{example}

\begin{example}[信号平移]
    取定时延b，$r(t)=e(t-b)$.
\end{example}
这也是时滞系统的例子，只不过不涉及微分的运算。

\begin{example}[傅里叶变换]
$r(t)=\mathcal{F} e(t)$.

在一些光学仪器中可以实现空间上的傅里叶变换，这时变量t是多维的空间变量，例如矩孔夫琅禾费衍射装置。
\end{example}

\subsection{系统的级联}\label{sub:cascade}
将一个系统的响应作为另一个系统的激励，这样的系统连接方式称为系统的\textbf{级联}(cascade)，下面就来研究这个话题。
\begin{theorem}[线性系统的级联]
    设有两线性系统K,L,将两系统级联之后，所得的新系统仍是线性系统。
\end{theorem}
\begin{proof}
    先利用$K$的线性性，再利用$L$的线性性，得到：
    \begin{align*}
    LK[ae_1+be_2]=L[aK[e_1]+bK[e_2]]=aLK[e_1]+bLK[e_2]
\end{align*}
\end{proof}

对核积分系统的级联也是对核积分系统，并且级联系统的核函数可以视作前级系统的核函数作为后级系统的激励信号输入后得到的响应函数：
\begin{proposition}[对核积分系统的级联]
    设有两对核积分系统K，L：
    \[K[e(t)](x)=\int_{a}^{b}k(x,y)e(y)\,dy,L[e(t)](x)=\int_{c}^{d}l(x,y)e(y)\,dy\]
则级联系统$H=LK$也是对核积分系统，并且核函数为\begin{equation*}
    h(x,z)=L_y(k(y,z))(x)=\int_{c}^{d}l(x,y)k(y,z)\,dy
\end{equation*}
其中$L_y$表示$k(x,y)$作为y的函数输入系统L，输出为x,y的双变量函数。
\end{proposition}

\begin{proof}\begin{align*}
    H[e(z)](x)&=LK[e(z)](x)\\
    &=\int_{c}^{d}l(x,y)\left(\int_{a}^{b}k(y,z)e(z)\,dz\right)\,dy\\
    &=\int_{a}^{b}e(z)\,dz\int_{c}^{d}l(x,y)k(y,z)\,dy\\
    &=\int_{a}^{b}h(x,z)e(z)\,dz
\end{align*}
\end{proof}

函数与$\delta$的卷积还是它本身，这就好比信号在进入任何一个系统前，先经过了一个与$\delta$卷积的系统，利用这一观点即可得到\textbf{叠加定理}(superposition thoerem)\cite{brad_osgood_ee261}，由于它是泛函分析中对应定理的特殊形式，有时也将它称为\textbf{施瓦兹核定理}(Schwartz kernel thoerem)。
\begin{theorem}[叠加定理/施瓦兹核定理]
    \begin{equation*}
        r(x)=\int_{-\infty}^{\infty}k(x,y)e(y)\,dy
    \end{equation*}
其中$k(x,y)$是线性系统H的核函数。
\end{theorem}
\begin{proof}\begin{align*}
    H[e(y)]&=H[e*\delta(y)]\\
    &=H\left[\int_{-\infty}^{\infty}\delta(x-y)e(y)\,dy\right]\\
    &=\int_{-\infty}^{\infty}H[\delta(x-y)]e(y)\,dy
\end{align*}
令$k(x,y)=H[\delta(x-y)]$即得。
\end{proof}

特别地，有限区间上的对核积分系统也满足叠加定理，此时核函数满足：
\[k(x,y)=k(x,y)\chi_{[a,b]}(y)\]

由于对核积分系统是线性的，系统的线性性就等价于系统能够用对核积分来描述，并且核函数总为
$$k(x,y)=H_x[\delta(x-y)]$$

\begin{proposition}[线性系统的等价定义]
    “系统$H$是线性系统”等价于“响应是激励对核$k(x,y)$积分”。
\end{proposition}

对于时不变系统，这个定理具有更加优美的形式。将系统视为线性系统，其核函数为
\[k(x,y)=H_x[\delta(x-y)]=H[\tau_y\delta(x)]=\tau_y H[\delta(x)]=h(x-y)\]
其中$h$表示系统的单位脉冲响应。带入叠加定理，即得：
\begin{align*}
    r(x)&=\int_{-\infty}^{\infty}h(x-y)e(y)\,dy\\
    &=(h*e)(t)
\end{align*}

也就是说，\textbf{信号经过线性时不变系统相当于与这个系统的单位脉冲响应做卷积}，这样，对线性时不变系统的研究归结为对单位脉冲响应的研究，极大地降低了研究的难度。另一方面，卷积描述的系统也是线性时不变的：\begin{align*}
    H[\tau_b e(y)](x)=(h*\tau_b e)(x)=\tau_b(h*e)(x)=\tau_b H[e(y)]
\end{align*}
而线性性是显然的。总之，我们有：
\begin{proposition}[线性时不变系统的等价定义]
    “系统$H$是线性时不变系统”等价于“响应是激励与单位脉冲响应$h$的卷积”，即$r(t)=(h*e)(t)$。
\end{proposition}

线性时不变系统的单位阶跃响应是单位脉冲响应的积分，可以利用此关系式进行单位脉冲响应或单位阶跃响应的求解，下一小节中将给出一个利用此方法的例子。

\begin{proposition}[单位阶跃响应]
线性时不变系统的单位阶跃响应为
$$g(t)=\int_{-\infty}^{t}h(y)\,dy$$
其中$h(t)$为系统的单位脉冲响应。
\end{proposition}
\begin{proof}
\begin{align*}
    g(t)&=H[u(t)]=(u*h)(t)\\
    &=\int_{-\infty}^{\infty}u(t-y)h(y)\,dy\\
    &=\int_{-\infty}^{t}h(y)\,dy
\end{align*}
\end{proof}

最后，我们来说明系统的一个重要特性：\textbf{因果性}，并利用$r(t)=(h*e)(t)$给出线性时不变系统的因果性条件。

\begin{definition}[因果性]
一个系统在有激励时，才会出现响应，或者说$r(t_0)$仅依赖于$e(t)$在$t<t_0$时的值（这里不等号不取等是标准的定义方式，我们马上会看到它的作用），即
\begin{equation}
    \forall t_0,\left(e_1(t)=e_2(t)\ for\ t<t_0\right)\implies \left(r_1(t)=r_2(t)\ for\ t<t_0\right)\label{eq:4.1}
\end{equation}
这个条件称为\textbf{因果性条件}(casuality condition)。
\end{definition}

因果性表明，系统的响应不能先于激励出现，且仅仅依赖于先前的激励，从物理上讲，这样的系统才可能是可实现的。对于特殊的系统类型，因果性条件具有更加简洁的形式：

\begin{proposition}
对于线性系统，因果性条件等价于
\begin{equation}
    \forall t_0,\left(e(t)=0\ for\ t<t_0\right)\implies \left(r(t)=0\ for\ t<t_0\right)\label{eq:4.2}
\end{equation}
\end{proposition}
\begin{proof}
只要取$e_1(t)=e(t),e_2(t)=0$，就从条件\ref{eq:4.1}推出条件\ref{eq:4.2}；反过来，如果条件\ref{eq:4.2}成立，则任意$e_1,e_2$满足
$$e_1(t)=e_2(t)\ for\ t<t_0,\text{即}e_1(t)-e_2(t)=0\ for\ t<t_0$$
根据条件\ref{eq:4.2}，有
$$r_1(t)-r_2(t)=0\ for\ t<t_0,\text{即}r_1(t)=r_2(t)\ for\ t<t_0$$
\end{proof}

\begin{proposition}
对于线性时不变系统，因果性条件等价于
\begin{equation}
    \left(e(t)=0\ for\ t<0\right)\implies \left(r(t)=0\ for\ t<0\right)\label{eq:4.3}
\end{equation}
或者更简便的
\begin{equation}
    h(t)=0\ for\ t<0\label{eq:4.4}
\end{equation}
其中h为单位脉冲响应。
\end{proposition}
\begin{proof}
只要取$t_0=0$，就从条件\ref{eq:4.2}推出条件\ref{eq:4.3}；用时不变性，条件\ref{eq:4.3}可以推出条件\ref{eq:4.2}。

对于单位脉冲响应，
$$\delta=0\ for\ t<0\implies h(t)=0\ for\ t<0$$
反过来，如果$h(t)=0,e(t)=0\ for\ t<0$,则根据\ref{sub:intuition_of_convolution}中介绍的卷积的“支集相加”性质，有$r(t)=(h*e)(t)=0\ for\ t<0$.
\end{proof}
